\section{A pointed presentation and its Gram matrix}

Fix $x_0\in P$. Deleting the row $x_0$ from $A$ and appending the all-ones
block row gives a full-row-rank $v\times b$ matrix $\Cx$; the basis
$\{e_x-e_{x_0}\}\cup\{k e_{x_0}\}$ of $\Hk$ turns each incidence column into its
coordinates away from $x_0$ followed by a final $1$.

\begin{proposition}\label{prop:gp}
$\Delt(D)\cong\coker(\Cx)$, and different choices of $x_0$ differ by unimodular
row operations.

The Gram matrix $G=\Cx\Cx^{T}$ depends only on the ordinary
$2$-design parameters,
\[
G=\begin{pmatrix}\alpha I_{n}+\lambda J_{n} & r\,\mathbf 1_n\\
r\,\mathbf 1_n^{T} & b\end{pmatrix},\qquad
\det G=\alpha^{\,n}\,r/k,
\]
so $|\Delt(D)|^{2}$ divides $\alpha^{\,n}r/k$, and for every prime $p$,
$2\,\ell_p(\Delt)\le n\,v_p(\alpha)+v_p(r)-v_p(k)$.

Writing $R=\im(\Cx^{T})$ and
$\bar R=(R\otimes\QQ)\cap\ZZ^{b}$, the group $\bar R/R$ is canonically
Pontryagin dual to $\Delt(D)$ and embeds in $\coker(G)\cong R^{\#}/R$ as a
totally isotropic subgroup for the bilinear discriminant form
$\beta_G(z,w)=z^{T}G^{-1}w\bmod\ZZ$. \textup(No quadratic refinement is asserted
without parity hypotheses.\textup)
\end{proposition}

\begin{proof}
The pointed basis identifies $\Hk/L$ with $\coker(\Cx)$; the Gram computation is
the standard $2$-design identity, and $\det G=\alpha^{n}r/k$ follows from
$(rk-\lambda)b-nr^{2}=\alpha r/k$.

Since $\Cx$ has full row rank $v$, its finite cokernel has order $|\Delt|$ equal
to the greatest common divisor of the $v\times v$ minors of $\Cx$ (the top
determinantal divisor $d_v$). By the Cauchy--Binet formula, $\det
G=\det(\Cx\Cx^{T})$ equals the sum of the squares of those maximal minors; each
minor is divisible by $|\Delt|$, so every summand is divisible by $|\Delt|^{2}$,
whence $|\Delt|^{2}\mid\det G$. Taking $p$-adic valuations gives the primewise
bound $2\ell_p(\Delt)\le v_p(\det G)=n\,v_p(\alpha)+v_p(r)-v_p(k)$.

For completeness, the duality pairing is explicitly
\[
(\bar R/R)\times(\ZZ^v/\Cx\ZZ^b)\longrightarrow\QQ/\ZZ,
\qquad (\Cx^Tq+R,z+\Cx\ZZ^b)\longmapsto q^Tz\bmod\ZZ.
\]
The vector $q\in\QQ^v$ is unique, since $\Cx$ has full row rank.
Changing either representative changes the value by an integer. Smith normal
form shows that this pairing is perfect, hence identifies $\bar R/R$ with
$\Delt^\vee$.

Finally $\bar R\subseteq\ZZ^b$ has integral dot products,
so $\bar R\subseteq R^\#$ and the induced discriminant pairing vanishes on
$\bar R/R$. This is the bilinear saturation construction underlying the
overlattice framework~\cite{Nikulin}; no even-lattice assumption is used.
\end{proof}

Thus the abstract annihilation of Theorem~\ref{thm:hs} becomes a computable
structural invariant: a finite bilinear form determined by the parameters, with
the dual defect sitting inside it isotropically.
